\chapter{Complexity — Meyer Machine}
%%% generated by the blueprint pipeline from TCSlib.Complexity.CircuitComplexity.MeyerMachine
\section{Overview}

Declarations of the Lean module \texttt{TCSlib.Complexity.CircuitComplexity.MeyerMachine}.

\section{Declarations}

\begin{definition}[Table selection labels]
\lean{Complexity.Meyer.TabSel}
\label{Complexity.Meyer.TabSel}
\leanok
\uses{Complexity.TimeHierarchy.OutReg}
For a type $S$ and a natural number $k$, a table selection label is either a
state-and-output pair consisting of an optional state from $S$ and an output-register
value, or a symbol label consisting of an optional index in $\{0,\ldots,k-1\}$ and two
Boolean values.
\end{definition}

\begin{definition}[Symbol bit test]
\lean{Complexity.Meyer.symBit}
\label{Complexity.Meyer.symBit}
\leanok
Given a Boolean value indicating whether the symbol is blank and an optional Boolean
symbol value $v$, the symbol bit test is true exactly when $v$ is absent if the blank
indicator is true, and exactly when $v$ equals true if the blank indicator is false.
\end{definition}

\begin{definition}[Selection answer]
\lean{Complexity.Meyer.selAnswer}
\label{Complexity.Meyer.selAnswer}
\leanok
\uses{Complexity.Meyer.TabSel, Complexity.Meyer.symBit, Complexity.TimeHierarchy.OutReg}
Given a table selection label, an optional state, an output-register value, an optional
Boolean symbol, and a function assigning an optional Boolean symbol to each of the $k$
positions, the selection answer is determined as follows. For a state-and-output label
it is true exactly when both supplied values match the label; for a symbol label without
an index it is the symbol bit test on the supplied symbol; and for a symbol label with
an index it is the symbol bit test on the symbol at that position.
\end{definition}

\begin{definition}[Walk direction]
\lean{Complexity.Meyer.walkDir}
\label{Complexity.Meyer.walkDir}
\leanok
For a Boolean value `neg`, the walk direction is negative if `neg` is true and positive
otherwise.
\end{definition}

\begin{definition}[Selection table size]
\lean{Complexity.Meyer.nbitsOf}
\label{Complexity.Meyer.nbitsOf}
\leanok
\uses{Complexity.Meyer.TabSel}
For a finite type $S$ with decidable equality and a natural number $k$, the number of
bits associated with $S$ and $k$ is the cardinality of the set of table selection labels
for $S$ and $k$.
\end{definition}

\begin{theorem}[Positive number of selector bits]
\lean{Complexity.Meyer.nbitsOf_pos}
\label{Complexity.Meyer.nbitsOf_pos}
\leanok
\uses{Complexity.Meyer.nbitsOf}
\difficulty{1}
For every finite type $S$ with decidable equality and every natural number $k$, the
number of selector states for $S$ and $k$ is positive.
\end{theorem}

\begin{definition}[One-hot selector encoding]
\lean{Complexity.Meyer.encodeSel}
\label{Complexity.Meyer.encodeSel}
\leanok
\uses{Complexity.Meyer.TabSel, Complexity.Meyer.nbitsOf}
Let $S$ be a finite type and let $k$ be a natural number. The encoding of a selector
from $\mathrm{TabSel}(S,k)$ is the Boolean vector indexed by the
$|\mathrm{TabSel}(S,k)|$ positions that is true exactly at the position corresponding to
that selector and false elsewhere.
\end{definition}

\begin{theorem}[Injectivity of the selector encoding]
\lean{Complexity.Meyer.encodeSel_injective}
\label{Complexity.Meyer.encodeSel_injective}
\leanok
\uses{Complexity.Meyer.TabSel, Complexity.Meyer.encodeSel}
\difficulty{3}
Let $S$ be a finite type with decidable equality, let $k$ be a natural number, and let
the selector set consist of the two forms specified by $	exttt{TabSel}$: a state option
paired with an output-register value, or a tape-index option paired with two bits.
Encode each selector as the Boolean function on the selector set’s finite cardinality
that is $1$ at its own index under a fixed enumeration and $0$ elsewhere. This encoding
is injective.
\end{theorem}

\begin{definition}[Selector decoding]
\lean{Complexity.Meyer.decodeSel}
\label{Complexity.Meyer.decodeSel}
\leanok
\uses{Complexity.Meyer.TabSel, Complexity.Meyer.encodeSel, Complexity.Meyer.nbitsOf}
Let $S$ be a finite type with decidable equality, let $k$ be a natural number, and let
$b$ be a Boolean vector indexed by the $|\mathrm{TabSel}(S,k)|$ positions. If $b$ is the
encoding of some selector, decode it as such a selector; otherwise return the selector
of the form $\mathrm{sym}(\mathrm{none},\mathrm{false},\mathrm{false})$.
\end{definition}

\begin{theorem}[Selector encoding and decoding]
\lean{Complexity.Meyer.decodeSel_encodeSel}
\label{Complexity.Meyer.decodeSel_encodeSel}
\leanok
\uses{Complexity.Meyer.TabSel, Complexity.Meyer.decodeSel, Complexity.Meyer.encodeSel, Complexity.Meyer.encodeSel_injective}
\difficulty{3}
For a finite state type $S$ with decidable equality and any number $k$ of work tapes,
encoding a selector as its Boolean indicator vector and then decoding that vector
returns the original selector.
\end{theorem}

\begin{definition}[Machine transition states]
\lean{Complexity.Meyer.TState}
\label{Complexity.Meyer.TState}
\leanok
\uses{Complexity.TimeHierarchy.OutReg}
For a type $S$ and a natural number $n$, define the transition-state type to have the
following constructors: a selector state carrying an index in $\{0,\ldots,n-1\}$ and an
$n$-bit Boolean vector; an $f_A$ state carrying a Boolean and such a vector; an $f_B$
state carrying two Booleans and such a vector; copy, rewrite-$X$, rewrite-$C$, and
rewrite-output states, each carrying such a vector; decode and return states, each
carrying a vector, an element of $S$, an output-register value, and a Boolean; and
weak-decode and weak-return states, each carrying a vector, an optional element of $S$,
an output-register value, and a Boolean.
\end{definition}

\begin{definition}[Bit count of the machine state encoding]
\lean{Complexity.Meyer.nbits}
\label{Complexity.Meyer.nbits}
\leanok
\uses{Complexity.Meyer.nbitsOf}
For a finite-tape machine $M$ with Boolean alphabet, define its bit count to be the
number of bits specified by the state type of $M$ and its number of work tapes.
\end{definition}

\begin{definition}[Selector type for a finite-tape machine]
\lean{Complexity.Meyer.Sel}
\label{Complexity.Meyer.Sel}
\leanok
\uses{Complexity.Meyer.TabSel}
For a finite-tape machine $M$ over the Boolean alphabet, let its selector type be the
type of tabular selectors on the machine’s state space and with parameter $k$.
\end{definition}

\begin{definition}[Machine transition-state type]
\lean{Complexity.Meyer.St}
\label{Complexity.Meyer.St}
\leanok
\uses{Complexity.Meyer.TState, Complexity.Meyer.nbits}
For a Boolean finite-tape machine $M$, define its transition-state type to be the
transition-state type with parameters $S=M.\mathrm{State}$ and $n$ equal to its selector
bit count.
\end{definition}

\begin{definition}[Action on Extended Work Tapes]
\lean{Complexity.Meyer.tact}
\label{Complexity.Meyer.tact}
\leanok
\uses{Complexity.Meyer.St, Turing.Action}
Let $M$ be a finite-tape Turing machine over the Boolean alphabet, and let its auxiliary
control-state set be the state set used by the construction. Given an input-head
movement, an action for each of the $M.k$ work tapes, three further work-tape actions,
an optional output symbol, and an optional successor auxiliary control state, form the
action on $M.k+3$ work tapes that uses those work-tape actions in that order and retains
the supplied input-head movement, output symbol, and successor state.
\end{definition}

\begin{definition}[Idle work-tape actions]
\lean{Complexity.Meyer.idle}
\label{Complexity.Meyer.idle}
\leanok
For a natural number $n$, define the idle action on $n$ work tapes to assign to every
tape the pair consisting of no write instruction and zero head movement.
\end{definition}

\begin{definition}[Single-cell tape action]
\lean{Complexity.Meyer.one}
\label{Complexity.Meyer.one}
\leanok
\uses{Complexity.Meyer.idle}
For $n\in\mathbb{N}$, a position $j\in\{0,\ldots,n-1\}$, and a tape action $a$, define
an action on an $n$-cell array that assigns $a$ at position $j$ and leaves every other
position idle.
\end{definition}

\begin{definition}[Two-cell tape action]
\lean{Complexity.Meyer.two}
\label{Complexity.Meyer.two}
\leanok
\uses{Complexity.Meyer.one}
For $n\in\mathbb{N}$, positions $j,j'\in\{0,\ldots,n-1\}$, and tape actions $a,a'$,
define an action on an $n$-cell array by assigning $a$ at $j$ and then assigning $a'$ at
$j'$, leaving all other cells idle. If $j=j'$, the assignment at $j'$ determines the
resulting action at that position.
\end{definition}

\begin{definition}[Rejecting machine action]
\lean{Complexity.Meyer.reject}
\label{Complexity.Meyer.reject}
\leanok
\uses{Complexity.Meyer.St, Complexity.Meyer.idle, Complexity.Meyer.tact, Turing.Action}
For a finite Turing machine $M$ with Boolean tape symbols, define a one-step action for
a machine with three additional work tapes that leaves the input head and all work-tape
heads stationary, writes nothing on the work tapes, emits the output symbol false, and
halts.
\end{definition}

\begin{definition}[Transition of the encoded simulator]
\lean{Complexity.Meyer.tabTr}
\label{Complexity.Meyer.tabTr}
\leanok
\uses{Complexity.Meyer.St, Complexity.Meyer.decodeSel, Complexity.Meyer.idle, Complexity.Meyer.nbits, Complexity.Meyer.one, Complexity.Meyer.reject, Complexity.Meyer.selAnswer, Complexity.Meyer.tact, Complexity.Meyer.two, Complexity.Meyer.walkDir, Complexity.TimeHierarchy.OutReg.push, Turing.Action, Turing.FinTM.virtualMove, Turing.FinTM.virtualNextTag}
Let $M$ be a finite Turing machine with Boolean tape symbols, and let $q$ be a state of
its encoded simulator. Given an optional Boolean input symbol and the symbols scanned on
its $k+3$ work tapes, the definition assigns the simulator’s one-step action, including
its tape movements and writes, output, and successor state or halting behavior.
\end{definition}

\begin{definition}[Transition table machine]
\lean{Complexity.Meyer.tabTM}
\label{Complexity.Meyer.tabTM}
\leanok
\uses{Complexity.Meyer.St, Complexity.Meyer.nbitsOf_pos, Complexity.Meyer.tabTr, Turing.FinTM}
Given a finite Turing machine $M$ with Boolean tape symbols, define a finite Turing
machine with $M.k+3$ work tapes and state set $\text{St}(M)$. Its initial state is the
selector state indexed by $0$ with all the associated bits false, and its transition
function is the transition rule $\text{tabTr}(M)$.
\end{definition}

\begin{definition}[Simulator configuration]
\lean{Complexity.Meyer.tcfg}
\label{Complexity.Meyer.tcfg}
\leanok
\uses{Complexity.Meyer.St, Complexity.Meyer.tabTM, Turing.Cfg}
Let $M$ be a finite Turing machine with Boolean symbols, and let $y$ be a finite Boolean
input. A configuration of the simulator machine is specified by a state of its simulator
state set or no state, an input-head position among the $|y|+2$ available positions, the
contents and head positions of each of the $M.k$ simulated work tapes and three
auxiliary tapes, and a finite Boolean output string; the simulated tapes precede the
auxiliary tapes in the simulator's work-tape tuple.
\end{definition}

\begin{definition}[Conditional tape write]
\lean{Complexity.Meyer.wr}
\label{Complexity.Meyer.wr}
\leanok
For an optional optional Boolean symbol $w$, a tape contents function
$\tau:\mathbb{Z}\to\mathrm{Option}(\mathrm{Bool})$, and a position $z\in\mathbb{Z}$,
define the updated tape to be unchanged if $w$ is absent; if $w$ is present, replace the
contents at $z$ by its inner optional symbol and leave all other positions unchanged.
\end{definition}

\begin{theorem}[Effect of a Tabulation Action]
\lean{Complexity.Meyer.tact_apply}
\label{Complexity.Meyer.tact_apply}
\leanok
\uses{Complexity.Meyer.St, Complexity.Meyer.tact, Complexity.Meyer.tcfg, Complexity.Meyer.wr, Turing.Action.apply, Turing.moveInputPos}
\difficulty{5}
Let $M$ be a finite Turing machine with Boolean symbols, and let $y$ be a finite Boolean
input word. Choose an input-head movement, a write instruction and head movement for
each of the $M.k$ work tapes and three auxiliary tapes, an optional output bit, and an
optional successor state; also choose a configuration with state, input-head position,
tape contents and head positions, and output word. Applying the action formed from these
instructions to this configuration yields the configuration with the chosen successor
state, input head moved by the clamped input-head movement, each original and auxiliary
tape updated at its head exactly when a write is prescribed, each head moved by its
prescribed amount, and the optional output bit appended to the output word.
\end{theorem}

\begin{theorem}[Tabulation machine step]
\lean{Complexity.Meyer.tabTM_step}
\label{Complexity.Meyer.tabTM_step}
\leanok
\uses{Complexity.Meyer.St, Complexity.Meyer.tabTM, Complexity.Meyer.tabTr, Complexity.Meyer.tcfg, Turing.Action.apply, Turing.Cfg.inputSymbol, Turing.Cfg.workTapeSymbols, Turing.MultiTapeTM.step}
\difficulty{4}
Let $M$ be a finite Turing machine with Boolean symbols, and let $y$ be a finite Boolean
input word. For any state $q$ of the tabulation machine, valid input-head position $p$,
contents and head positions for the work tapes of $M$, contents and head positions for
three auxiliary tapes, and Boolean output word, one step of the tabulation machine from
the resulting configuration is the configuration obtained by applying the action
prescribed by its transition function to that configuration. The transition function is
evaluated at the symbol scanned on the input tape and the symbols scanned on the
original and auxiliary work tapes.
\end{theorem}

\begin{theorem}[Input Symbol Independent of Work Tapes]
\lean{Complexity.Meyer.tcfg_inputSymbol}
\label{Complexity.Meyer.tcfg_inputSymbol}
\leanok
\uses{Complexity.Meyer.St, Complexity.Meyer.tcfg, Turing.Cfg.inputSymbol}
\difficulty{1}
For a Boolean finite Turing machine $M$, any Boolean input word, state, valid input-head
position, original work-tape contents and head positions, three auxiliary work-tape
contents and head positions, and output word, the symbol read under the input head in
the associated configuration is unchanged when all work tapes and the output are
replaced by blank tapes and an empty output, keeping the same input-head position.
\end{theorem}
